\title{Direct Preference Optimization: Your Language Model is Secretly a Reward Model}

\begin{abstract}
Although large-scale unsupervised language models (LMs) acquire extensive world knowledge and demonstrate certain reasoning abilities, exerting precise influence over their outputs remains challenging due to the fully unsupervised approach used in their training. Current approaches for enhancing such controllability typically involve collecting human-annotated comparisons of model outputs and adjusting the unsupervised LM to better match these preferences, commonly utilizing reinforcement learning from human feedback (RLHF). Nevertheless, RLHF introduces considerable complexity and instability: it requires first constructing a reward model aligned with human judgments, and then applying reinforcement learning to tune the LM so as to maximize this reward, all while ensuring minimal deviation from the base model. In this work, we propose a novel way of parameterizing the reward model in RLHF, which makes it possible to recover the associated optimal policy analytically. This reformulation enables solving the standard RLHF problem via a straightforward classification objective. The approach we introduce, termed \textit{Direct Preference Optimization} (DPO), is robust, effective, and computationally efficient, as it removes the necessity for language model sampling during fine-tuning and reduces the need for extensive hyperparameter search. Empirically, we demonstrate that DPO is able to adapt LMs to human preferences as well as, or better than, state-of-the-art techniques. Importantly, DPO outperforms RLHF with PPO in tasks involving sentiment control, and yields comparable or superior performance for summarization and single-turn dialogue, all while offering substantial simplification in both implementation and training.
\end{abstract}

\section{Introduction}
Unsupervised language models (LMs) developed on immense corpora have demonstrated unexpectedly broad competencies~\citep{chowdhery2022palm, brown2020language, touvron2023llama,bubeck2023sparks}. Nevertheless, their training material—generated by people with differing motives, backgrounds, and expertise—often contains behaviors we may not wish an AI system to replicate. For instance, in the case of an AI code assistant, it is desirable that it can \textit{identify} frequent coding errors to provide appropriate corrections, but during code generation, our aim is for it to emulate the exemplary programming skills, even if such skills are rare within the dataset. Similarly, we would want a language model to be \textit{informed} about a misconception prevalent among 50\% of individuals, but clearly, we do not want it to assert this mistaken belief as truth in half of its responses. Hence, it is essential to distinguish and select a model's \emph{preferred actions and answers} from its much broader pool of \textit{learned knowledge and potential outputs}, thereby ensuring the construction of AI systems that are safe, effective, and easily directed \citep{ouyang2022training}. Although current methodologies typically employ reinforcement learning (RL) to align language models with human preferences, we demonstrate that the prevailing RL-based preference objective can, in fact, be realized precisely using a straightforward binary cross-entropy formulation, thereby streamlining the overall preference learning process.

Broadly, the standard approach to aligning model outputs involves using datasets of human-curated preferences to encourage desirable behaviors. This process—known as preference learning—supplements the initial phase of extensive, unsupervised pre-training conducted on massive text datasets. While a naïve solution to preference learning might involve supervised tuning based on high-quality human response demonstrations, reinforcement learning from human or AI feedback (RLHF/RLAIF; \citep{christiano2017deep,bai2022constitutional}) has emerged as the most effective paradigm. Within RLHF, a reward model is trained to reflect human judgments, and reinforcement learning is subsequently applied to refine the model to favor responses rated highly by this reward model, while regularizing to maintain proximity to the pre-trained model. Despite delivering impressive results in both dialogue and coding applications, RLHF introduces substantial complexity beyond supervised fine-tuning—it requires training several models and repeated sampling from model outputs during optimization, resulting in high computational demands.

In this work, we present a method to train language models directly according to human preferences, circumventing the need for explicit reward models or reinforcement learning approaches. Our method, termed \textit{Direct Preference Optimization} (DPO), is designed to implicitly optimize the same underlying objective found in standard RLHF algorithms—namely, maximizing reward while enforcing a KL-divergence constraint—yet it is notably simpler to both implement and train. The essence of DPO lies in its update rule: it enhances the relative log-likelihood of preferred responses over less favored ones. Importantly, this update is weighted dynamically on a per-example basis, mitigating model collapse commonly observed with a naive probability ratio formulation. DPO, similar to prevailing approaches, is grounded in a theoretical preference framework, such as the Bradley-Terry model \cite{bradley1952rankanalysis}, to gauge the alignment between a given reward function and observed human preference data. Unlike standard methods—which employ the preference model to create a loss for reward model training, followed by policy training to optimize this learned reward—DPO re-parameterizes the loss in terms of the policy itself. Given a collection of preference-labeled data, DPO employs a straightforward binary cross entropy loss to train the policy, effectively yielding an optimal policy for an implicit reward inferred from the preference labels.

The principal innovation introduced in this paper is Direct Preference Optimization (DPO), a streamlined, RL-free procedure for preference-based language model training. Through empirical studies, we demonstrate that DPO matches or surpasses the effectiveness of contemporary techniques, including those based on PPO-driven RLHF, across a range of applications such as sentiment adjustment, summarization, and conversational agents, using models of up to 6B parameters.

Language models trained with self-supervised objectives and scaled up in size have demonstrated the ability to tackle some tasks with no additional training data (zero-shot; \citep{radford2019language}) or by conditioning on a few demonstration examples (few-shot; \citep{gpt3,megatron,chowdhery2022palm}). Nevertheless, their effectiveness on downstream applications and their alignment with human intent can be markedly enhanced through fine-tuning on curated datasets composed of explicit instructions paired with human-generated responses \citep{mishra-etal-2022-cross,sanh2022multitask,chung2022scaling,thoppilan2022lamda}. This process, commonly referred to as `instruction-tuning,' allows large language models (LLMs) to handle instructions that were not present during fine-tuning and generally improves user experience \citep{chung2022scaling}. 

Despite the positive outcomes of instruction tuning, collecting \textit{relative} human assessments of answer quality tends to be more feasible than gathering expert-labeled demonstrations. Consequently, later research has leveraged datasets containing human preference comparisons to further fine-tune LLMs, leading to advancements in tasks such as translation \citep{kreutzer-etal-2018-reliability}, summarization \citep{stiennon2022learning,ziegler2020finetuning}, story generation \citep{ziegler2020finetuning}, and instruction following \citep{ouyang2022training,ramamurthy2023is}. The prevailing strategy involves first training a neural reward model to fit these preferences—typically through a preference aggregation method like the Bradley-Terry model \citep{bradley1952rankanalysis}—and then adjusting the language model itself to maximize the learned reward using reinforcement learning techniques. Such techniques often include REINFORCE \citep{williams1992reinforce}, proximal policy optimization (PPO; \cite{schulman2017proximal}), or adaptations thereof \citep{ramamurthy2023is}. 

A related avenue exploits LLMs, themselves already aligned with instructions and human feedback, to produce synthetic preference datasets oriented towards specific goals such as safety or harmlessness \citep{bai2022constitutional}; in this case, weak human supervision, typically in the form of rubric-based annotation guidelines, is employed. These approaches bridge two key strands of research: the development of reinforcement learning algorithms for fine-tuning language models across diverse objectives \citep{Ranzato2015SequenceLT,paulus2018a,wu2018learning}, and general frameworks for learning from human preference data \citep{christiano2017deep,kupcsik2018learning}. While leveraging relative human preferences is promising, the process of refining large language models using reinforcement learning poses substantial practical obstacles. This work introduces a theoretically principled alternative for optimizing models with respect to relative preferences, circumventing the need for reinforcement learning.

Beyond applications involving language, learning policies from preferences has been explored extensively in both bandit and reinforcement learning paradigms, leading to a variety of proposed methods. Contextual bandit frameworks that leverage preferences or rankings—rather than direct rewards—are categorized as contextual dueling bandits (CDB; \cite{yue2012karmed,dudik2015contextual}). In these cases, where absolute reward signals are unavailable, theoretical characterizations of optimality in CDBs rely on the concept of the \textit{von Neumann winner}: a policy that guarantees an expected win probability of at least 50\% against any alternative policy \citep{dudik2015contextual}. One distinction of the CDB regime is that feedback in the form of preference labels arrives sequentially in an online manner. In contrast, typical learning-from-human-preferences scenarios involve training on a static, finite collection of action pairs annotated with offline human preferences \citep{yan2022human}. Relatedly, in \textit{preference-based RL} (PbRL), learning occurs from binary preferences derived from an \textit{unknown} underlying scoring function instead of direct reward feedback \citep{BusaFekete2014,ruiz2023dueling}. A range of PbRL algorithms has been developed, with several approaches accommodating off-policy preference datasets. These commonly entail a two-stage process of estimating the latent scoring function (serving as a reward model) followed by policy optimization with respect to it \citep{jain2013learning,BusaFekete2014,christiano2017deep,sadigh2017active,kupcsik2018learning}. In contrast, we introduce a unified procedure that eschews explicit reward modeling, and instead directly optimizes the policy to conform to the provided preferences.

\section{Preliminaries}\label{section:prelims}

We provide an overview of the RLHF pipeline as described by \citeauthor{ziegler2020finetuning} and subsequently adopted in works such as \citep{stiennon2022learning, bai2022training, ouyang2022training}. This pipeline generally consists of three key steps: 1) supervised fine-tuning (SFT), 2) the gathering of preference data coupled with reward model training, and 3) reinforcement learning-based optimization.

\textbf{SFT}: The RLHF process usually commences with the supervised fine-tuning of a pre-existing language model using curated datasets tailored to downstream applications (e.g., dialogue, summarization), yielding a model denoted by $\pi^\text{SFT}$.

\textbf{Reward Modelling Phase}: In the subsequent stage, prompts $x$ are input to the SFT model, which generates answer pairs $(y_1, y_2)\sim \pi^\text{SFT}(y \mid x)$. These responses are shown to human annotators, who are tasked with indicating their preference (notated as $y_w\succ y_l \mid x$, where $y_w$ and $y_l$ denote, respectively, the preferred and less-preferred responses from $(y_1, y_2)$). These preference judgments are assumed to be governed by an unobserved reward function $r^*(y, x)$, which is not directly accessible. Several strategies have been proposed for capturing such preferences, with the Bradley-Terry (BT) model \cite{bradley1952rankanalysis} being a prominent example. Other ranking models such as Plackett-Luce \citep{plackett1975analysis, luce2012individual} can also be employed if sets of ranked answers are available. Within the BT framework, the distribution over human preferences $p^*$ is represented as:

Given a fixed dataset of pairwise comparisons $\mathcal{D}=\bigl\{x^{(i)}, y_w^{(i)}, y_l^{(i)}\bigr\}_{i=1}^N$ drawn from $p^*$, we can introduce a parameterized reward function $r_{\phi}(x, y)$ and fit its parameters by maximizing the likelihood of the observed data. Casting the problem as binary classification leads to a negative log-likelihood objective:

\begin{equation}\label{eq:reward_model}
    \mathcal{L}_R(r_{\phi}, \mathcal{D}) = -\mathbb{E}_{(x, y_w, y_l)\sim \mathcal{D}}\bigl[\log \sigma(r_{\phi}(x, y_w)- r_{\phi}(x, y_l))\bigr]
\end{equation}

where $\sigma$ denotes the logistic sigmoid function. Within large language model settings, $r_{\phi}(x, y)$ is generally initialized from a SFT model $\pi^\text{SFT}(y \mid x)$, augmented by an additional linear transformation atop the final transformer output to yield a scalar reward prediction \cite{ziegler2020finetuning}. Previous work normalizes reward predictions to ensure $\mathbb{E}_{x,y\sim \mathcal{D}}\left[r_\phi(x, y)\right] = 0$ over all $x$, thereby reducing reward function variance.

\textbf{RL Fine-Tuning Phase}: In the subsequent RL phase, the acquired reward model furnishes supervisory signals to further tune the language model. Following established approaches~\citep{jaques2017sequence, jaques2020human}, we optimize

\begin{equation}\label{eq:RL}
\max_{\pi_{\theta}}  \mathbb{E}_{x\sim \mathcal{D}, y\sim \pi_{\theta}(y \mid x)}\bigl[r_{\phi}(x, y)\bigr] - \beta\mathbb{D}_{\textrm{KL}}\bigl[\pi_{\theta}(y\mid x)\mid \mid \pi_\text{ref}(y\mid x)\bigr],
\end{equation}

where $\beta$ modulates how far $\pi_\theta$ can diverge from a reference policy $\pi_\text{ref}$, typically set to the initial SFT model $\pi^\text{SFT}$. Commonly, $\pi_\theta$ is initialized as $\pi^\text{SFT}$. The KL term is crucial for restricting updates so that the policy stays close to the regions where the reward model remains reliable and also preserves diversity in the generations, thereby guarding against mode collapse onto a small set of high-reward outputs. Since the sampling process in language models is inherently discrete and the objective is non-differentiable, reinforcement learning is required for optimization. The prevailing methodology \citep{ziegler2020finetuning, stiennon2022learning, bai2022training, ouyang2022training} constructs a composite reward of the form ${r(x, y) = r_{\phi}(x, y) -\beta (\log \pi_{\theta}(y\mid x) - \log \pi_\text{ref}(y\mid x))}$ and maximizes it with Proximal Policy Optimization (PPO) \cite{schulman2017proximal}.

\section{Direct Preference Optimization}\label{sec:DPO}

Recognizing the limitations of RL-based algorithms for large-scale tasks such as language model fine-tuning, we seek a more straightforward method for optimizing policies based directly on preferences. In contrast to standard RLHF pipelines—where one first learns a reward function and then maximizes it via reinforcement learning—our approach hinges on a specific parameterization of the reward model that allows for analytical computation of its optimal policy, entirely sidestepping iterative RL updates. As detailed below, the core innovation is to employ an exact relationship that links reward models to their corresponding optimal policies, thereby facilitating a change-of-variables that transforms a reward-based loss into a policy-based loss. This technique obviates the need to train a standalone reward network, while retaining the capacity to model human preferences under frameworks like the Bradley-Terry model. Consequently, the policy network simultaneously serves as both the generative model and an implicit estimator of reward.

\textbf{DPO Objective Derivation.} We begin with the conventional RL objective outlined in Eq.~\ref{eq:RL}, assuming a generic reward function $r$. Building upon earlier findings~\citep{peters2007reinforcement, peng2019advantage, korbak2022reinforcement, go2023aligning}, it can be shown that the solution which optimally satisfies the KL-regularized reward maximization objective in Eq.~\ref{eq:RL} adopts the form:

\begin{equation}\label{eq:op_policy}
    \pi_r(y\mid x) = \frac{1}{Z(x)}\pi_\text{ref}(y\mid x)\exp\left(\frac{1}{\beta}r(x, y)\right),
\end{equation}

where the normalization constant $Z(x) =\sum_{y}\pi_\text{ref}(y\mid x)\exp\left(\frac{1}{\beta}r(x, y)\right)$ is often called the partition function. Full details of this derivation are available in Appendix \ref{app:derivation1}. When substituting the MLE estimate $r_{\phi}$ for the true reward $r^*$, estimating $Z(x)$ remains computationally intensive~\citep{korbak2022reinforcement, go2023aligning}, making practical application of this policy challenging. To mitigate this, one can rewrite Eq.~\ref{eq:op_policy}, re-expressing the reward function using its corresponding optimal policy $\pi_r$, the reference policy $\pi_\text{ref}$, and the (generally intractable) partition function $Z(\cdot)$. Taking the logarithm of both sides of Eq.~\ref{eq:op_policy} yields, after algebraic rearrangement:

\begin{equation}\label{eq:main_eq}
    r(x,y) =\beta \log \frac{\pi_r(y\mid x)}{\pi_\text{ref}(y\mid x)} + \beta \log Z(x).
\end{equation}

This form can be directly applied to the true reward $r^*$ and its associated optimal policy $\pi^*$. Notably, the Bradley-Terry model depends solely on differences in reward for two possible completions; that is, ${p^*(y_1 \succ y_2 \mid x) = \sigma(r^*(x, y_1) - r^*(x, y_2))}$. By plugging the representation from Eq.~\ref{eq:main_eq} for $r^*(x,y)$ into the preference function in Eq.~\ref{eq:bradley-terry}, the terms involving the partition function cancel, enabling an expression for human preference likelihoods solely via the optimal policy $\pi^*$ and the reference policy $\pi_\text{ref}$. Therefore, under the Bradley-Terry formulation, the ideal RLHF policy $\pi^*$ is described by:

\begin{equation}\label{eq:objective}
    p^*(y_1\succ y_2 \mid x)=\frac{1}{1 + \exp\left(\beta \log \frac{\pi^*(y_2\mid x)}{\pi_\text{ref}(y_2\mid x)} - \beta \log \frac{\pi^*(y_1\mid x)}{\pi_\text{ref}(y_1\mid x)}\right)}
\end{equation}

A detailed derivation is found in Appendix~\ref{app:derivation2}. Although Eq.~\ref{eq:objective} relies on the Bradley-Terry model, an analogous derivation applies under broader Plackett-Luce models~\citep{plackett1975analysis, luce2012individual}, as detailed in Appendix~\ref{app:plackett_luce_models}.

Expressing the probability of human preferences as a function of the optimal policy (rather than the reward), we can now define a maximum likelihood objective tailored for a parameterized policy $\pi_\theta$. Mirroring the logic of the reward modeling paradigm (see Eq.~\ref{eq:reward_model}), we derive the policy training objective as:

\begin{equation}\label{eq:optimum_model}
    \mathcal{L}_\text{DPO}(\pi_{\theta}; \pi_\text{ref}) = -\mathbb{E}_{(x, y_w, y_l)\sim \mathcal{D}}\left[\log \sigma \left(\beta \log \frac{\pi_{\theta}(y_w\mid x)}{\pi_\text{ref}(y_w\mid x)} - \beta \log

In this formulation, we learn an implicit reward through an alternative parameterization, where the optimal policy directly corresponds to $\pi_\theta$. Additionally, our approach is mathematically equivalent to training a reparameterized Bradley-Terry model, which yields advantageous theoretical guarantees—such as consistency under appropriate preference data distributions \cite{bong2022generalized}. Further discussion of the theoretical properties of DPO and its connections to related literature is provided in Section~\ref{sec:theory}.

\textbf{What effect does the DPO update have?} To gain insight into the mechanics of DPO, we consider the gradient of the loss $\mathcal{L}_\text{DPO}$. Explicitly, the gradient with respect to $\theta$ is given by:

\begin{multline*}\label{eq:gradient}
    \nabla_\theta \mathcal{L}_\text{DPO}(\pi_\theta;\pi_\text{ref}) = \\ -\beta\mathbb{E}_{(x, y_w, y_l) \sim \mathcal{D}} \bigg[\underbrace{\sigma(\hat{r}_\theta(x, y_l) - \hat{r}_\theta (x, y_w))}_\text{greater emphasis when reward assignment is inaccurate}\bigg[\underbrace{\nabla_\theta\log \pi(y_w \mid x)}_\text{boost probability of $y_w$} - \underbrace{\nabla_\theta\log\pi(y_l \mid x)}_\text{reduce probability of $y_l$}\bigg]\bigg],
\end{multline*}

with $\hat{r}_\theta(x, y) = \beta \log \frac{\pi_\theta(y \mid x)}{\pi_\text{ref}(y \mid x)}$ serving as the implicitly defined reward from the language model $\pi_\theta$ and the reference $\pi_\text{ref}$ (elaborated in Section~\ref{sec:theory}). Conceptually, this gradient update encourages the model to make the preferred continuations $y_w$ more probable and penalizes the less favored completions $y_l$. Crucially, the strength of each training example is modulated by the degree to which the model's current reward assignment overestimates $y_l$ relative to $y_w$, controlled by $\beta$—effectively prioritizing examples where the model most misranks the completions, while also factoring in the KL constraint. Empirically, our results indicate that this form of weighting is essential: omitting it—yielding a simpler, unweighted update—leads to degeneration in language model outputs (see Appendix Table~\ref{tab:unlikelihood_generations}).

\textbf{DPO summary.} The standard DPO workflow involves: 1) Generating response samples $y_1, y_2 \sim \pi_\text{ref}(\cdot \mid x)$ for each prompt $x$, then annotating these samples using human preferences to create an offline preference dataset $\mathcal{D} = \{x^{(i)}, y_w^{(i)}, y_l)^{(i)}\}_{i=1}^N$; and 2) updating the parameters of the language model $\pi_\theta$ by minimizing $\mathcal{L}_\text{DPO}$ with respect to the given $\pi_\text{ref}$, $\mathcal{D}$, and a specified $\beta$. 
In applied settings, practitioners often wish to leverage existing public preference datasets instead of collecting new samples and human feedback. Since these datasets are typically constructed with completions from $\pi^\text{SFT}$, the standard approach is to set $\pi_\text{ref} = \pi^\text{SFT}$ if such a model exists. Conversely, in situations where $\pi^\text{SFT}$ is unavailable, $\pi_\text{ref}$ is obtained by maximizing the likelihood over the preferred completions ${(x, y_w)}$ in the dataset, i.e., ${\pi_\text{ref} = \argmax_{\pi}\mathbb{E}_{x, y_w \sim \mathcal{D}}\left[\log \pi(y_w \mid x)\right]}$. This initialization strategy aims to reduce the discrepancy between the actual, but inaccessible, reference distribution and the $\pi_\text{ref}$ employed by DPO. Readers can find a comprehensive discussion of implementation specifics and hyperparameter settings in Appendix~\ref{app:implementation}.

\section{Theoretical Analysis of DPO} In what follows, we provide a deeper interpretation of the DPO approach, offer a theoretical foundation for its design, and demonstrate how DPO addresses certain limitations present in actor-critic methods commonly used for RLHF, such as PPO~\cite{schulman2017proximal}.

\label{sec:theory}

\subsection{Your Language Model Is Secretly a Reward Model} DPO unifies reward learning and policy optimization into a single maximum likelihood estimation framework, obviating the need for explicit reward modeling and separate RL training. Notably, the DPO objective defined in Eq.~\ref{eq:main_eq} corresponds mathematically to a Bradley-Terry model, where rewards are parameterized as $r^*(x, y) = \beta \log\frac{\pi^*_\theta(y \mid x)}{\pi_\text{ref}(y \mid x)}$. Training involves directly updating the parametric language model $\pi_{\theta}$, which is equivalent to the process of reward model optimization described in Eq.~\ref{eq:reward_model} after reparameterization. This section develops the theoretical rationale behind this change of variables, establishes that the resulting formulation imposes no additional restrictions on the set of learnable reward models, and demonstrates that this framework allows recovery of the optimal policy exactly. We begin by formalizing an equivalence relation on reward functions.

\begin{definition}
We define two reward functions $r(x, y)$ and $r'(x, y)$ as equivalent if and only if there exists a function $f$ such that ${r(x, y)-r'(x, y) = f(x)}$.
\end{definition}

This definition indeed forms an equivalence relation, thereby dividing the set of reward functions into distinct equivalence classes. The following lemmas hold:

\begin{lemma}\label{lemma:same_prefrence} Within the framework of Plackett-Luce—and in particular, the Bradley-Terry preference model—any two reward functions that belong to the same equivalence class generate identical preference distributions.
\end{lemma}

\begin{lemma}\label{lemma:same_policy}
    Any pair of reward functions from the same equivalence class will yield the same optimal policy for the constrained RL problem.
\end{lemma}

These results follow straightforwardly, and we provide detailed proofs in Appendix \ref{app:lemma1}. The first lemma highlights the well-recognized issue of under-specification present in the Plackett-Luce family \cite{plackett1975analysis}. To address this ambiguity, it is standard practice to incorporate extra identifiability constraints to ensure meaningful MLE estimates in Eq. \ref{eq:reward_model} \cite{bong2022generalized}. The second lemma implies that all functions within the same equivalence class lead to an identical optimal policy; thus, for our main objective, it suffices to recover any representative reward function from the correct equivalence class. We establish the following theorem in Appendix~\ref{app:thm1}:

\begin{theorem}\label{thm:main}
    Provided mild assumptions hold, every reward equivalence class compatible with the Plackett-Luce model (including Bradley-Terry as a special case) can be characterized through the reparameterization ${r(x, y) = \beta \log \frac{\pi(y\mid x)}{\pi_\text{ref}(y\mid x)}}$ for some model $\pi(y\mid x)$ and a fixed reference model $\pi_\text{ref}(y \mid x)$.
\end{theorem}

\begin{sproof}
    Let $r(x, y)$ be any reward function, which defines an associated optimal model $\pi_r(y \mid x)$ according to Eq. \ref{eq:op_policy}. We now demonstrate that a reward function from the equivalence class of $r$ admits the specified reparameterization. Specifically, let us introduce the projection $f$ as 
\begin{equation}
    f(r; \pi_\text{ref}, \beta)(x, y) = r(x, y) - \beta\log\sum_{y}\pi_\text{ref}(y\mid x)\exp\left(\frac{1}{\beta}r(x, y)\right)
\end{equation}
This operator $f$ effectively normalizes the reward function using the log-partition function of $\pi_r$. Since this normalization term depends solely on the context $x$, the function $f(r; \pi_\text{ref}, \beta)(x, y)$ remains within the equivalence class of $r(x, y)$. By substituting the right-hand side of Eq.~\ref{eq:main_eq} (which universally applies to all reward functions), we see that $f(r; \pi_\text{ref}, \beta)(x, y) = \beta \log \frac{\pi_r(y\mid x)}{\pi_\text{ref}(y\mid x)}$. Consequently, this projection constructs a reward function within the equivalence class of $r$ having the desired reparameterized form, thereby confirming the generality of our reward model parameterization.
\end{sproof}

An alternative interpretation of Theorem~\ref{thm:main} is that it precisely identifies which reward function, among those in a given equivalence class, is chosen by the DPO reparameterization. Specifically, it selects the reward function that satisfies

\begin{equation}\label{eq:lag_p}
     \sum_{y}\underbrace{\pi_\text{ref}(y\mid x)\exp\left(\frac{1}{\beta}r(x, y)\right)}_{=\pi(y\mid x)\text{, using Thm.~\ref{thm:main} reparam.}} = 1,
\end{equation}

which ensures that $\pi(y\mid x)$ is indeed a well-defined probability distribution: it is non-negative and sums to one. Observing Eq.~\ref{eq:op_policy}, it is clear that Eq.~\ref{eq:lag_p} corresponds to the normalization constant (partition function) for the optimal policy defined by the reward function $r(x, y)$. The fundamental insight underpinning the DPO approach is that by imposing suitable constraints on the generally under-determined Plackett-Luce (and in particular, Bradley-Terry) class of preference models, we can both retain the richness of possible reward models and explicitly ensure the normalization constant in Eq.~\ref{eq:op_policy} is analytically computable for every prompt $x$.

\subsection{Instability of Actor-Critic Algorithms} Our framework can also illuminate sources of instability inherent in standard actor-critic algorithms, such as PPO, used for RLHF. Adhering to the RLHF methodology, we concentrate on the RL fine-tuning phase described in Section~\ref{section:prelims}. This perspective can be linked to the control as inference paradigm~\cite{levine2018reinforcement}, which addresses the constrained RL setting outlined in~\ref{eq:RL}. We parameterize our model as $\pi_{\theta}(y\mid x)$, aiming to minimize $\mathbb{D}_{\text{KL}}[\pi_{\theta}(y|x) \mid \mid \pi^*(y\mid x)]$, where $\pi^*$ denotes the optimal policy induced by $r_{\phi}(y, x)$ per Eq.~\ref{eq:optimum_model}. After appropriate algebraic manipulation, the resulting objective can be written as:

\begin{equation}\label{eq:AC}
    \max_{\pi_{\theta}}\mathbb{E}_{\pi_{\theta}(y\mid x)}\left[\underbrace{r_{\phi}(x, y) -\beta\log\sum_{y}\pi_\text{ref}(y\mid x)\exp\left(\frac{1}{\beta}r_{\phi}(x, y)\right)}_{f(r_{\phi}, \pi_\text{ref}, \beta)} - \underbrace{\beta\log\frac{\pi_{\theta}(y\mid x)}{\pi_\text{ref}(y\mid x)}}_{\text{KL}}\right]
\end{equation}

The objective considered here aligns with that optimized in previous studies 
\citep{ziegler2020finetuning, stiennon2022learning, bai2022training, ouyang2022training}, making use of the DPO-equivalent reward defined over the reward class $r_{\phi}$. Within this context, the normalization component in $f(r_{\phi}, \pi_\text{ref}, \beta)$ can be viewed as representing the soft value function corresponding to the reference policy $\pi_\text{ref}$. Although the normalization term does not modify the set of optimal solutions, its omission may result in the policy gradient having large variance, potentially leading to instability in the training process. One option is to learn a value function to account for this normalization; however, this approach often introduces optimization challenges. Alternatively, earlier research has addressed normalization by referencing a human completion baseline, which amounts to using a single-sample Monte-Carlo estimator for the normalization. The DPO reparameterization, on the other hand, produces a reward formulation that eliminates the need for such baselines.

\section{Experiments}
This section presents an empirical investigation into the effectiveness of DPO in training policies directly from human preferences. Initially, we focus on a controlled text-generation environment to compare how well DPO balances reward maximization and KL-divergence minimization relative to widely used preference learning techniques like PPO. Subsequently, we extend our assessment to encompass larger models and more challenging RLHF scenarios, such as tasks involving summarization and dialogue. Our experiments indicate that, even with minimal hyperparameter adjustment, DPO frequently matches or surpasses established baselines, including RLHF implemented with PPO and selection of the best out of $N$ sampled trajectories under a learned reward function. Prior to presenting these empirical findings, we detail our experimental methodology; supplementary information is available in Appendix~\ref{app:exp_details}.

\textbf{Tasks.} We conduct experiments on three open-ended text generation tasks. Across all tasks, the algorithms are trained to optimize a policy using a dataset of preferences, denoted as $\mathcal{D}=\bigl\{x^{(i)}, y_w^{(i)}, y_l^{(i)}\bigr\}_{i=1}^N$. In the case of \textbf{controlled sentiment generation}, each $x$ corresponds to the beginning of a movie review sourced from the IMDb dataset \cite{maas-EtAl:2011:ACL-HLT2011}, and the objective is for the policy to produce an output $y$ exhibiting positive sentiment. For robust evaluation, preference pairs are constructed by generating outputs and using a pre-trained sentiment classifier, ensuring that $p(\text{positive}\mid x,y_w)>p(\text{positive}\mid x,y_l)$. Supervised fine-tuning (SFT) is performed on GPT-2-large, training to convergence using reviews from the IMDb train set (additional details are provided in App~\ref{app:sentiment_details}). In the \textbf{summarization} task, $x$ refers to a Reddit forum post, with the goal of the policy being to generate a summary $y$ that captures the post’s main ideas. We use the Reddit TL;DR summarization dataset \citep{volske-etal-2017-tl} and a collection of human preference data assembled by \citeauthor{stiennon2022learning}. For SFT, we rely on a model trained on human-authored summaries of forum posts\footnote{\url{https://huggingface.co/CarperAI/openai_summarize_tldr_sft}}, leveraging the TRLX \citep{leandro_von_werra_2023_7790115} library for reinforcement learning from human feedback (RLHF). The set of human preferences utilized originates from \citeauthor{stiennon2022learning}, but is based on samples from another SFT model that underwent similar training. Lastly, for \textbf{single-turn dialogue}, each $x$ is a user’s query, which may vary from technical questions on astrophysics to personal advice requests. Here, the policy must formulate an informative and engaging reply $y$. Experiments use the Anthropic Helpful and Harmless dialogue dataset \citep{bai2022training}, which includes 170k dialogues between a user and a virtual assistant. Each dialogue concludes with two assistant responses generated by a large, albeit unspecified, language model, accompanied by a label identifying the response preferred by a human annotator. As no pre-existing SFT model is provided for this task, we build one by fine-tuning an available language model exclusively on preferred responses.

\textbf{Evaluation.} Our experimental evaluation utilizes two complementary methodologies. For the controlled sentiment generation task, where the underlying reward function (provided by a sentiment classifier) is fully accessible, we assess each algorithm by charting the frontier of its achieved reward against the KL-divergence relative to the reference policy. This explicit frontier is feasible because we possess the exact form of the reward function. Conversely, in scenarios more akin to real-world applications—where the true reward function is not available—we assess each algorithm based on its \textit{win rate} over a baseline policy. In these cases, GPT-4 is used as a stand-in for human judgment to evaluate both summary quality (summarization task) and response helpfulness (dialogue task). Specifically, in the summarization task, the test set's reference summaries are employed as baselines; for dialogue, the dataset's preferred responses serve this role. Prior work indicates that language models can outperform traditional metrics as automated evaluators \citep{Chen2023ExploringTU}. Nevertheless, to substantiate our choice of GPT-4 for evaluation, we supplement our study with a human evaluation in Sec.~\ref{sec:human-judgments}. Our findings show that human judgments are highly consistent with those produced by GPT-4, with agreement rates that are comparable to or exceed the level of agreement among human annotators.

\textbf{Methods.} Alongside DPO, we benchmark a variety of established strategies for training language models according to human preferences. For baseline prompting approaches, we utilize zero-shot prompts with \textbf{GPT-J} \citep{gpt-j} for summarization and employ 2-shot prompts using \textbf{Pythia-2.8B} \citep{biderman2023pythia} in the dialogue experiments. We further include the \textbf{SFT} model, as well as \textbf{Preferred-FT}—a model refined via supervised fine-tuning on selected completions $y_w$, drawn from either the SFT outputs (for controlled sentiment and summarization) or a standard language model (for single-turn dialogue). Another method under evaluation is \textbf{Unlikelihood}~\citep{welleck2019neural}, which trains the policy to assign high likelihood to $y_w$ and concurrently reduce the probability of $y_l$, governed by a tunable coefficient $\alpha\in[0,1]$ that modulates the strength of the unlikelihood term. We also examine \textbf{PPO} \citep{schulman2017proximal}, which leverages a reward model trained on preference data, as well as \textbf{PPO-GT}, an oracle approach utilizing the true reward function available in the sentiment control experiments. Within these sentiment experiments, we implement PPO-GT both as an off-the-shelf solution \cite{leandro_von_werra_2023_7790115} and in an adapted variant that includes reward normalization and refined hyperparameters for optimized results; analogous modifications are applied when using PPO with learned rewards. To round out our comparison, we use the \textbf{Best of $N$} baseline, which involves sampling $N$ outputs from the SFT model (or Preferred-FT in the dialogue setting) and returning the candidate with the highest reward, as determined by a reward function learned from preferences. Although this method can yield strong results by dissociating reward modeling from PPO optimization, it is computationally intensive—even for moderate values of $N$—due to the necessity of generating $N$ responses per prompt at evaluation time.

\subsection{How well can DPO optimize the RLHF objective?}

The typical RLHF framework employs a KL-regularized reward maximization objective that aims to maximize the reward while simultaneously ensuring that the policy remains close to the reference policy. Consequently, fair comparisons across different algorithms must jointly assess both the achieved reward and the degree of policy shift as measured by KL divergence; simply attaining a marginally higher reward with a substantially increased KL is not necessarily preferable. In Figure~\ref{fig:frontier-tldr-main}, we present the trade-off frontier between reward and KL for a range of algorithms within the sentiment analysis context. Each algorithm is subjected to multiple independent training runs, each corresponding to a distinct setting of policy conservativeness—specifically, target KL values of $\{3,6,9,12\}$ for PPO, $\beta$ values in $\{0.05,0.1,1,5\}$, $\alpha$ in $\{0.05,0.1,0.5,1\}$ for unlikelihood, and varying random seeds for preferred-FT—amounting to a total of 22 separate runs. At intervals of 100 training steps up to convergence, we assess each policy against a suite of held-out prompts, computing both the mean reward according to the ground-truth reward function and the average sequence-level KL\footnote{This refers to the cumulative sum of KL-divergences across timesteps in a sequence.} relative to the reference policy, $\text{KL}\left(\pi\mid \mid \pi_\text{ref}\right)$. Our analysis reveals that DPO consistently constructs the most favorable reward-KL frontier, delivering maximal reward with minimal KL. This outcome stands out for several reasons: Firstly, DPO and PPO formally target the same underlying optimization problem, yet DPO displays markedly superior efficiency—DPO's position along the reward/KL spectrum uniformly surpasses that of PPO. Secondly, DPO achieves an improved frontier compared to PPO even in scenarios where PPO is granted access to the ground-truth rewards (PPO-GT).

\subsection{Can DPO scale to real preference datasets?}
\label{sec:dpo-real-datasets}

We proceed by assessing how well DPO fine-tunes language models on tasks of summarization and single-turn dialogue. In the context of summarization, standard automatic evaluation metrics like ROUGE often have limited alignment with human preference judgments~\citep{stiennon2022learning}, and previous research has shown that employing PPO for preference-based fine-tuning leads to higher-quality summaries. Our evaluation compares multiple approaches by generating completions on the test partition of the TL;DR summarization dataset and measuring the mean win rate relative to test set reference completions. For each method, completions are generated with sampling temperatures ranging from 0.0 to 1.0, and corresponding win rates are visualized in Figure~\ref{fig:frontier-tldr-main} (right). DPO, PPO, and Preferred-FT all leverage fine-tuning on the same GPT-J SFT checkpoint\footnote{\url{https://huggingface.co/CarperAI/openai_summarize_tldr_sft}}. Our results indicate that DPO achieves an approximate 61\% win rate at temperature 0.0, surpassing PPO’s peak win rate of around 57\% at its optimal temperature setting. Additionally, DPO outperforms the $N$ baseline’s maximum win rate. It is important to highlight that we did not conduct comprehensive tuning of the DPO $\beta$ hyperparameter, suggesting these findings may not fully represent DPO’s capabilities. We also observe that DPO is substantially less sensitive to temperature changes compared to PPO, whose performance deteriorates to baseline GPT-J levels at higher temperatures. Preferred-FT yields minimal gains over the initial SFT model. Furthermore, direct human evaluation of DPO versus PPO is discussed in Section~\ref{sec:human-judgments}; in this comparison, samples from DPO at temperature 0.25 are favored 58\% of the time over PPO samples at the same temperature.

For single-turn dialogue evaluation, we assess each method on the subset of the Anthropic HH dataset’s test split \citep{bai2022training} involving a single human-assistant exchange. In our GPT-4-based evaluations, the win rate for each approach is determined by comparing outputs to the test set's preferred completions. Due to the absence of an established SFT benchmark for this task, we initialize from a pre-trained Pythia-2.8B model. We employ Preferred-FT to fine-tune a reference model on selected completions to ensure alignment between the distribution of completions and model outputs, then proceed with DPO training. Additionally, we benchmark against the optimal selection from 128 Preferred-FT completions (with the Best of $N$ metric saturating at $N=128$ for this evaluation; refer to Appendix Figure~\ref{fig:best-of-n}), and we also test a 2-shot prompt setup with the base Pythia-2.8B model. Across methods, DPO achieves equal or superior win rates at their respective optimal sampling temperatures. Moreover, we evaluate an RLHF model, trained via PPO on the Anthropic HH dataset \footnote{\url{https://huggingface.co/reciprocate/ppo_hh_pythia-6B}} and made available by a reputable repository \footnote{\url{https://github.com/CarperAI/trlx/tree/main/examples/hh}}; however, no tested combination of prompt or sampling temperature led this PPO model to outperform the original Pythia-2.8B. Drawing on insights from TL;DR results and the shared reward function between the two strategies, we treat Best of 128 as an approximate indicator of PPO-level performance. Collectively, DPO stands out as the only method that not only achieves computational efficiency but also surpasses the preferred completions baseline in the Anthropic HH dataset, offering comparable or enhanced results relative to the more resource-intensive Best of 128 baseline. Figure~\ref{fig:dialogue-main} further illustrates DPO’s rapid convergence to peak performance.

\subsection{Generalization to a new input distribution}

To further explore the robustness of PPO and DPO under distributional shifts, we test the policies trained on Reddit TL;DR summarization by applying them to a different dataset: news articles from the test split of the CNN/DailyMail dataset \citep{nallapati-etal-2016-abstractive}. The evaluation uses the top-performing sampling temperatures from TL;DR (0 and 0.25). Table~\ref{tab:ood} reports the findings. For assessment, we measured the GPT-4 win rate in relation to the reference summaries provided in the datasets, employing the same GPT-4 (C) prompt used in the Reddit TL;DR evaluation, with the term ``forum post'' switched to ``news article.'' DPO continues to substantially outperform PPO on this unseen data. These results serve as preliminary evidence that DPO policies maintain generalization ability comparable to PPO policies, despite DPO not making use of extra unlabeled Reddit TL;DR prompts as PPO does.

\subsection{Validating GPT-4 judgments with human judgments}
\label{sec:human-judgments}
We conduct a user study to assess the trustworthiness of GPT-4’s evaluation criteria, leveraging outcomes from the TL;DR summarization task and comparing two distinct GPT-4 prompts. The \textbf{GPT-4 (S)} (simple) prompt asks which summary most effectively conveys the key information from the post, while the \textbf{GPT-4 (C)} (concise) prompt in addition queries which summary is more succinct. We examine the latter prompt due to GPT-4’s inclination, when using the simple prompt, to favor longer and more repetitive summaries than humans typically do. The full wording of these prompts is available in Appendix~\ref{app:prompts}. We conduct three comparisons, featuring the top-performing system (DPO, temp. 0.25), the lowest-performing (PPO, temp. 1.0), and an intermediate model (SFT, temp. 0.25), each benchmarked against PPO with greedy sampling (its optimal temperature). This approach is intended to sample a broad range of output qualities; all three are pitted against the PPO baseline. Analysis indicates that, for both prompts, GPT-4’s preferences generally align with those of human raters to a degree comparable with human inter-rater agreement, supporting the use of GPT-4 as a surrogate for manual evaluation (since human annotators were limited, only DPO and PPO-1 were evaluated by multiple humans). On the whole, win rates computed using the \textbf{GPT-4 (C)} prompt more closely reflect human judgment, so this prompt is selected for the primary results in Section~\ref{sec:dpo-real-datasets}. Further methodological details regarding the human evaluation—including the rater interface and the list of human participants—are provided in Appendix~\ref{app:human-study}.

\section{Discussion}
Preference-based learning constitutes an effective and scalable approach for developing language models that are both powerful and aligned. In this work, we presented DPO, a streamlined framework for training language models via preference data without relying on reinforcement learning techniques. Instead of recasting preference learning as a reinforcement learning task to leverage existing RL algorithms, DPO establishes a correspondence between language model policies and reward functions, facilitating the direct optimization of models to conform to human preferences through a straightforward cross-entropy loss. This approach avoids reinforcement learning while preserving generality. Remarkably, DPO achieves performance on par with, or superior to, current RLHF methods—such as those utilizing PPO—with minimal hyperparameter adjustment, substantially lowering the barriers to training language models based on human feedback.

\textbf{Limitations \& Future Work.} Several significant avenues for further investigation arise from our findings. One question concerns the extent to which DPO-trained policies exhibit out-of-distribution generalization, particularly in comparison to methods that optimize explicit reward functions. Preliminary evidence indicates that DPO generalizes comparably to PPO-based approaches, but thorough evaluation is required. For instance, it remains to be seen whether leveraging self-labeling techniques in the DPO framework can enable effective utilization of unannotated prompts. Another aspect warranting further study is how over-optimization of reward surfaces in the direct preference optimization framework, and whether the modest performance dip observed in Figure~\ref{fig:dialogue-main}-right exemplifies this phenomenon. Although our experiments were limited to models with up to 6B parameters, scaling DPO to cutting-edge models with vastly more parameters represents a compelling future direction. Additionally, our results suggest that the prompt can influence the win rates reported by GPT-4; more research could illuminate optimal strategies for soliciting reliable automated judgments. Lastly, DPO’s potential extends to a variety of applications outside language modeling from human preferences, such as training generative models in diverse modalities.